\documentclass[conference]{IEEEtran}
\IEEEoverridecommandlockouts
\usepackage{cite}
\usepackage{amsmath,amssymb,amsfonts}
\usepackage{algorithmic}
\usepackage{graphicx}
\usepackage{textcomp}
\usepackage{xcolor}
\def\BibTeX{{\rm B\kern-.05em{\sc i\kern-.025em b}\kern-.08em
    T\kern-.1667em\lower.7ex\hbox{E}\kern-.125emX}}

\begin{document}

\title{Cloud Seal Enhanced: A Privacy-Preserving Multi-Tenant Cloud Deduplication Framework with AI-Driven Detection, Blockchain Auditability, and Post-Quantum Security}

\author{\IEEEauthorblockN{Vihangi Patterson}
\IEEEauthorblockA{\textit{Department of Computing} \\
\textit{Informatics Institute of Technology}\\
(University of Westminster, UK)\\
Colombo, Sri Lanka \\
vihango.20220011@iit.ac.lk}
\and
\IEEEauthorblockN{Pubudu De Silva}
\IEEEauthorblockA{\textit{Supervisor} \\
\textit{Informatics Institute of Technology}\\
(University of Westminster, UK)\\
Colombo, Sri Lanka \\
pubudu.s@iit.ac.lk}
}

\maketitle

\begin{abstract}
Cloud storage costs continue to escalate due to data redundancy, yet traditional deduplication techniques compromise privacy when encryption is applied. This paper presents Cloud Seal Enhanced, a novel multi-tenant deduplication framework that resolves the encryption-deduplication paradox through convergent encryption while integrating three advanced security mechanisms: AI-driven encrypted similarity detection, distributed blockchain audit trails, and NIST-approved post-quantum cryptography. Our system employs a two-layer Convolutional Neural Network for detecting encrypted file similarities beyond deterministic hashing, a Proof-of-Authority blockchain for immutable transaction logging, and CRYSTALS-Kyber-768 integration validated via simulation for quantum-resistant key encapsulation. Experimental validation demonstrates 42--55\% storage reduction with sub-millisecond duplicate detection via Bloom filters (1.2\% false positive rate at 10,000 items), 100\% blockchain tamper detection, zero premature deletions across multi-tenant scenarios, and a high throughput of 29,608 transactions per second (TPS). The lightweight CNN achieves an Area Under the Curve (AUC) of 0.797 and 100\% recall for near-duplicate identification. This work represents a novel integrated architecture combining these technologies for cloud storage security, providing both immediate cost savings and future-proof protection.
\end{abstract}

\begin{IEEEkeywords}
Cloud Deduplication, Convergent Encryption, Multi-Tenant Security, Post-Quantum Cryptography, Blockchain Audit Trail, Reference Counting, Kyber-768, Bloom Filter
\end{IEEEkeywords}

\section{Introduction}
\subsection{Motivation}
Global cloud storage demand reached 137.3 billion dollars in 2024 with projected growth to 376.4 billion dollars by 2030 \cite{r1}. Enterprise organizations face massive storage costs, with 40 to 80 percent of data being redundant duplicates \cite{r2}. Traditional deduplication techniques achieve significant storage reduction but fail when encryption is applied, creating the encryption-deduplication paradox: organizations must choose between security (encryption) and efficiency (deduplication).

Furthermore, the emergence of quantum computing poses an existential threat to current cryptographic systems. NIST standardized post-quantum algorithms in August 2024 \cite{r3}, yet few systems integrate these protections. The ``harvest now, decrypt later'' attack vector means adversaries can collect encrypted data today and decrypt it when quantum computers become available in 10 to 15 years \cite{r4}.

\subsection{Research Gap}
Current approaches to encrypted deduplication address isolated challenges but fail to provide holistic solutions \cite{r5}. A comprehensive analysis of recent literature reveals the following gaps:

\begin{itemize}
    \item \textbf{DupLESS} \cite{r6}: Implements convergent encryption with centralized key servers to prevent confirmation attacks. Li et al. \cite{r10} extended this with randomized convergent encryption to mitigate brute-force attacks. While effective for single-tenant scenarios, these architectures introduce a single point of failure and lack both quantum-resistant cryptography and tamper-evident audit mechanisms.
    \item \textbf{SALIGP} \cite{r7}: Employs genetic algorithms for AI-driven similarity detection combined with blockchain logging. The genetic programming approach exhibits exponential search spaces, making real-time deduplication impractical. Furthermore, it relies on classical encryption which offers limited post-quantum protection.
    \item \textbf{SeDaSC} \cite{r8}: Proposes a cryptographic server architecture for secure cross-tenant data sharing. The centralized trust model creates vulnerability and provides no quantum resistance.
    \item \textbf{Secure Dedup} \cite{r9}: Develops an ECC-based secure deduplication scheme. While eliminating centralized dependencies, elliptic curve cryptography remains vulnerable to Shor's algorithm, exposing archived data to future quantum decryption.
\end{itemize}

Puzio et al. \cite{r11} introduced proof-of-ownership protocols integrated with convergent encryption, while Stanek et al. \cite{r12} proposed threshold cryptography for distributed key management. Recent surveys confirm that few existing frameworks simultaneously integrate AI-based optimization, blockchain-based accountability, and post-quantum cryptographic protection for multi-tenant encrypted cloud deduplication.

\subsection{Contributions}
This paper makes the following novel contributions:
\begin{enumerate}
    \item \textbf{Integrated Architecture:} A novel framework combining AI-driven similarity detection, blockchain auditability, and post-quantum cryptography for encrypted cloud deduplication.
    \item \textbf{Multi-Tenant Reference Counting:} Ownership tracking mechanism that prevents premature deletion while enabling cross-tenant deduplication visibility.
    \item \textbf{Hybrid Encryption Scheme:} Backward-compatible integration of classical AES-256 and quantum-resistant Kyber-768.
    \item \textbf{Lightweight AI Detection:} A 2-layer Siamese CNN for encrypted similarity detection operating on 2048-dimensional statistical features, achieving 0.797 AUC without homomorphic encryption overhead.
\end{enumerate}

\section{System Architecture}
\subsection{Overview}
Cloud Seal Enhanced employs a microservices architecture comprising five core components: Encryption Engine, Deduplication Manager, Blockchain Ledger, AI Detection Module, and Storage Backend.

\begin{figure}[htbp]
\centerline{\includegraphics[width=\linewidth]{mm.png}}
\caption{Cloud Seal Enhanced High-Level System Architecture}
\label{fig:architecture}
\end{figure}

\subsection{Encryption Engine}
The Encryption Engine implements a hybrid scheme combining classical and post-quantum cryptography. For file uploads:
\begin{equation}
K_{conv} = \text{SHA-256}(T_{id} \parallel T_{sec} \parallel M)
\end{equation}
\begin{equation}
C = \text{AES-256-GCM}(M, K_{conv})
\end{equation}
where $M$ is plaintext content, $C$ is ciphertext, $T_{id}$ is the tenant identifier, and $T_{sec}$ is the tenant secret. This builds on the convergent encryption model introduced by Douceur et al. \cite{r13}.

When sharing between tenants, the system's Kyber-768 integration is validated via simulation for quantum-safe key exchange:
\begin{equation}
(C_{kyb}, K_{shrd}) = \text{Kyber-768.Encap}(PK_{recv})
\end{equation}
This ensures that even if a quantum adversary compromises routing, file decryption remains infeasible without breaking Kyber's MLWE-based security \cite{r14}.

\subsection{Deduplication Manager}
1) \textbf{Bloom Filter:} A probabilistic data structure \cite{r15} for efficient O(1) duplicate detection. For $n=10,000$ and $p=0.01$, $m=95,851$ bits and $k=7$ hash functions are utilized.

2) \textbf{Reference Counter:} Tracks ownership across tenants:
\begin{equation}
\text{RefCount}(CID) = |\{T_i : T_i \in \text{Owners}(CID)\}|
\end{equation}
Physical deletion occurs only when $\text{RefCount}(CID) = 0$.

3) \textbf{AI Similarity Detector:} Following the Siamese architecture introduced by Bromley et al. \cite{r16}, a lightweight 2-layer Siamese CNN maps 2048-dimensional Z-score normalized features to a 128-dimensional embedding space:
\begin{equation}
h = \max(0, W_{1} \cdot \phi(C) + b_{1})
\end{equation}
\begin{equation}
z = \frac{W_{2} \cdot h + b_{2}}{\|W_{2} \cdot h + b_{2}\|}
\end{equation}
\begin{equation}
\text{sim}(C_1, C_2) = z_1 \cdot z_2
\end{equation}
A cosine similarity threshold of 0.85 generates advisory warnings for near-duplicates, while exact deduplication requires strict cryptographic hash matching.

\subsection{Blockchain Ledger}
Implements Proof-of-Authority consensus. The block structure prevents tampering: Any modification to $\text{Block}_j$ changes $H_j$, invalidating all subsequent blocks.

\section{Implementation}
The system is implemented using Python 3.11, FastAPI, and cryptographic libraries. Upload workflow evaluates the Bloom Filter for probabilistically identifying duplications, then executes the CNN for near-duplicate advisory. 

\section{Experimental Evaluation}
\subsection{Experimental Setup}
Synthetic datasets of text, images, and binary files were generated with controlled mutations (reversed, shuffled, appended) to simulate near-duplicate attacks and normal redundancies. Tests were conducted on an Apple M2 with 16 GB memory.

\subsection{AI Detection Accuracy}
Evaluating the Siamese CNN on a dedicated verification dataset of 44 file pairs revealed an AUC (Area Under the Receiver Operating Characteristic Curve) of $0.797$. The AI correctly distinguished complete statistical differentiation with negative scores ranging from $0.07$ to $0.32$, while pure duplicates maintained similarities $>0.998$. Utilizing a 0.85 threshold ensures 100\% recall against near-duplicates on the dataset, offering strong advisory insight without compromising data integrity via auto-deduplication. 

\subsection{Threat Mitigation Validation}
Five critical threats were systematically tested:
\begin{itemize}
    \item \textbf{Confirmation Attack:} Tenant-salted keys ensured different organizations produced unique ciphertexts for identical plaintext.
    \item \textbf{Cross-Tenant Leakage:} Verification proved zero unauthorized decryption across varying keys.
    \item \textbf{Audit Log Tampering:} The blockchain component consistently detected 100\% of simulated block mutations.
    \item \textbf{Key Derivation Weakness:} The SHA-256 derivation yielded a 46.5\% cryptographically secure avalanche effect.
\end{itemize}

\subsection{Performance Efficiency}
For end-to-end deduplication checking, processing times were minimal: $0.11$ms for 1KB files and $1.95$ms for 1MB files. The Bloom Filter lookup operated uniformly at $0.001$ms. The Deduplication Manager achieved between 42\% and 55\% storage reduction across synthetic distributions, while the PoA blockchain ledger sustained a high throughput of 29,608 transactions per second (TPS), cementing its viability for enterprise-grade scalability.

\section{Discussion and Limitations}
While Cloud Seal Enhanced successfully demonstrated the feasibility of integrated PoA blockchains alongside AI-assisted deduplication, the system currently operates as a Proof-of-Concept.

\textbf{Limitations:} 
First, the AI components were evaluated on strictly synthesized datasets (44 pairs) limiting broad-spectrum real-world generalization; scaling this up to $>10,000$ distinct document clusters remains necessary. 
Second, IPFS content-addressing and Kyber-768 cryptographic operations were validated via simulation within a managed environment to verify structural APIs rather than actively deployed in distributed cluster conditions. Future production rollouts will necessitate leveraging hardware-optimized libraries (e.g., \textit{liboqs-python}) and true IPFS daemon integration to mitigate large-scale latency effects.

\section{Conclusion}
This paper presented Cloud Seal Enhanced, a novel framework combining AI similarity, blockchain auditability, and post-quantum cryptographic components for encrypted cloud deduplication. Experimental results demonstrate a 0.797 AUC for near-duplicate clustering, 42--55\% storage reduction, and 29,608 blockchain TPS. Ultimately, this framework ensures multi-tenant security while significantly reducing untracked data redundancy.

\begin{thebibliography}{00}
\bibitem{r1} Statista, ``Cloud Storage Market Size Worldwide from 2021 to 2030,'' Market Research Report, 2024.
\bibitem{r2} IDC, ``Enterprise Storage Efficiency and Deduplication Trends,'' White Paper, 2023.
\bibitem{r3} NIST, ``Post-Quantum Cryptography Standardization: Module-Lattice-Based Key-Encapsulation Mechanism Standard,'' FIPS 203, August 2024.
\bibitem{r4} M. Mosca and M. Piani, ``Quantum Threat Timeline Report,'' Global Risk Institute, 2018.
\bibitem{r5} X. Chen et al., ``A Survey on Encrypted Data Deduplication in Cloud Storage,'' IEEE Access, vol. 12, pp. 45123--45145, 2024.
\bibitem{r6} M. Bellare, S. Keelveedhi, and T. Ristenpart, ``DupLESS: Server-Aided Encryption for Deduplicated Storage,'' in Proc. 22nd USENIX Security Symposium, 2013, pp. 179--194.
\bibitem{r7} Y. Zhang, L. Wang, and X. Liu, ``Enhancing Cloud Security and Deduplication Efficiency with SALIGP and Cryptographic Authentication,'' Scientific Reports, vol. 15, article 1847, Jan 2025.
\bibitem{r8} A. Rodriguez, C. Martinez, and J. Garcia, ``SeDaSC: Secure Data Sharing with Cryptographic Server for Cloud Deduplication,'' Journal of Information Security and Applications, vol. 68, article 103518, Sept 2024.
\bibitem{r9} J. Liu, N. Asokan, and B. Pinkas, ``Secure Deduplication of Encrypted Data without Additional Independent Servers,'' in Proc. ACM SIGSAC Conference on Computer and Communications Security, 2015, pp. 874--885.
\bibitem{r10} J. Li, X. Chen, M. Li, J. Li, P. Lee, and W. Lou, ``Secure Deduplication with Efficient and Reliable Convergent Key Management,'' IEEE Transactions on Parallel and Distributed Systems, vol. 25, no. 6, pp. 1615--1625, June 2014.
\bibitem{r11} P. Puzio, R. Molva, M. Onen, and S. Loureiro, ``ClouDedup: Secure Deduplication with Encrypted Data for Cloud Storage,'' IEEE Transactions on Cloud Computing, vol. 8, no. 4, pp. 1165--1178, Oct 2020.
\bibitem{r12} J. Stanek, A. Sorniotti, E. Androulaki, and L. Kencl, ``A Secure Data Deduplication Scheme for Cloud Storage,'' in Proc. Financial Cryptography and Data Security, Springer, 2014, pp. 99--118.
\bibitem{r13} J.R. Douceur, A. Adya, W.J. Bolosky, D. Simon, and M. Theimer, ``Reclaiming Space from Duplicate Files in a Serverless Distributed File System,'' in Proc. 22nd International Conference on Distributed Computing Systems (ICDCS), 2002, pp. 617--624.
\bibitem{r14} G. Alagic et al., ``Status Report on the Third Round of the NIST Post-Quantum Cryptography Standardization Process,'' NIST Internal Report 8413, 2022.
\bibitem{r15} B.H. Bloom, ``Space/Time Trade-offs in Hash Coding with Allowable Errors,'' Communications of the ACM, vol. 13, no. 7, pp. 422--426, July 1970.
\bibitem{r16} J. Bromley, I. Guyon, Y. LeCun, E. S\"ackinger, and R. Shah, ``Signature Verification using a Siamese Time Delay Neural Network,'' in Advances in Neural Information Processing Systems (NeurIPS), vol. 6, 1993.
\end{thebibliography}
\vspace{12pt}
\end{document}
